% Part: first-order-logic
% Chapter: proof-systems
% Section: tableaux

\documentclass[../../../include/open-logic-section]{subfiles}

\begin{document}

\iftag{FOL}
      {\olfileid[zh]{fol}{prf}{tab}}
      {\olfileid[zh]{pl}{prf}{tab}}

\olsection{\usetoken{P}{tableau}}

虽然许多!!{derivation}系统处理的是!!{sentence}的排列，!!{tableau}处理的却是!!{signed formula}。!!^a{signed formula}是由一个真值符号（$\True$ 或 $\False$）和一个!!a{sentence}构成的二元组
\[
\sFmla{\True}{!A} \text{ 或 } \sFmla{\False}{!A}.
\]
!!^a{tableau}由排列在向下分支的树中的!!{signed formula}组成。它以若干\emph{假设}开始，接着通过应用某条推理规则，从其上方的某个!!{signed formula}得到新的!!{signed formula}。每条规则允许我们在一个支的末端添加一个或多个!!{signed formula}，或者并排添加两个!!{signed formula}——此时支一分为二，两个新建的!!{signed formula}构成两支的末端。

对某个复合!!{signed formula}应用一条规则，会添加其若干直接子!!{formula}对应的!!{signed formula}。它们成对出现，对两个符号各有一条规则。例如，$\TRule{\True}{\land}$ 规则作用于 $\sFmla{\True}{!A \land !B}$，允许在含 $\sFmla{\True}{!A \land !B}$ 的任意支的末端同时添加!!{signed formula} $\sFmla{\True}{!A}$ 和~$\sFmla{\True}{!B}$，而规则 $\TRule{\False}{!A \land !B}$ 则允许通过并排添加 $\sFmla{\False}{!A}$ 和 $\sFmla{\False}{!B}$ 来使支分裂。!!^a{tableau}是闭的，如果它的每一条支都含有匹配的一对!!{signed formula} $\sFmla{\True}{!A}$ 与 $\sFmla{\False}{!A}$。

基于!!{tableau}的 $\Proves$ 关系定义如下：$\Gamma \Proves !A$ 当且仅当存在某个有穷集~$\Gamma_0 = \{!B_1, \dots, !B_n\} \subseteq \Gamma$，使得存在一个闭的!!{tableau}用于如下假设
\[
\{\sFmla{\False}{!A}, \sFmla{\True}{!B_1}, \dots, \sFmla{\True}{!B_n}\}
\]
例如，下面是一个表明 $\Proves (!A \land !B) \lif !A$ 的闭!!{tableau}：
\begin{oltableau}
  [\sFmla{\False}{(\formula{A} \land \formula{B}) \lif \formula{A}}, just = \TAss
    [\sFmla{\True}{\formula{A} \land \formula{B}}, just = {\TRule{\False}{\lif}[1]}
      [\sFmla{\False}{\formula{A}}, just = {\TRule{\False}{\lif}[1]}
        [\sFmla{\True}{\formula{A}}, just = {\TRule{\True}{\lif}[2]}
          [\sFmla{\True}{\formula{B}}, just = {\TRule{\True}{\lif}[2]}, close]
        ]
      ]
    ]
  ]
\end{oltableau}

在!!{tableau}演算中，一个集合 $\Gamma$ 是不一致的，如果存在某个有穷子集 $\{!B_1, \dots, !B_n\} \subseteq \Gamma$，使得有一个用于假设
\[
\{\sFmla{\True}{!B_1}, \dots, \sFmla{\True}{!B_n}\}
\]
的闭!!{tableau}。

!!^{tableau}s 于 1950 年代由\zhFullName{Evert Beth}{埃弗特·贝特}{贝特}和\zhFullName{Jaakko Hintikka}{亚科·欣蒂卡}{欣蒂卡}各自独立发明，并由\zhFullName{Raymond Smullyan}{雷蒙德·斯穆里安}{斯穆里安}加以简化和推广。它们用起来非常容易，因为构造!!a{tableau}是一个非常系统的过程。由于!!{tableau}的系统性，它们也很适合用计算机实现。不过，!!a{tableau}往往难以阅读，它们与证明的联系有时也不易看出。这一方法也相当一般化，许多不同的逻辑都有!!{tableau}系统。!!^{tableau}s 还帮助我们找到满足给定!!{sentence}（或语句集）的!!{structure}：如果该集合是可满足的，它就不会有闭的!!{tableau}，即任何!!{tableau}都会有一条开放的支。满足条件的!!{structure}可以被从开放的支「读取」，只要该支上所有可能应用的规则都已经应用。表列演算与矢列演算之间也有非常密切的联系：本质上，一个闭的!!{tableau}是矢列演算中的一个简略的!!{derivation}，只是倒过来书写。

\end{document}
